\section{Implementation and reproducibility}
\paragraph{Sampling individual training queries.}
Borg individual-model queries are built hourly, sampling up to 2,000 active
tasks without replacement from each state, with inverse inclusion weights.
Sampling uses a fixed seed and never uses future labels. Complete-population
scenario evaluation uses every active task at six-hour origins. BPI queries
contain every unresolved case at daily origins. Consequently, the individual
development table and the complete-population scenario table are different
estimands; the latter is not obtained by scoring a small sampled population.
The bounded development adapters produce 1,833,102 Borg queries and 519,142
BPI queries across training and development periods. The pre-final Borg prefix
contains 115,400,646 task-event rows; the BPI prefix contains 976,753 events.

\paragraph{Marginal training.}
HGB uses 250 boosting iterations, at most 31 leaves, a minimum of 100 samples
per leaf, learning rate 0.05, and L2 regularization 1.0, with early stopping
disabled. The recent Borg feature extension counts nine event types over
1/5/15/60-minute windows for personal, global, job, and machine scopes, using
log-one-plus counts. Temporal encoders use width 64, two Mixer blocks or a GRU,
typed event categories, continuous event features, observed time lags, and
padding masks. Histories are ordered oldest to newest, capped at 16 events per
scope. Parameters use AdamW with learning rate 0.001, weight decay $10^{-4}$,
batch size 2,048, gradient clipping at 5, and at most 30 epochs. The checkpoint
minimizes integrated Brier in an earlier calibration window, with patience 5.
The same window fits the optional scalar temperature in log bounds
$[-1.5,1.5]$. Final forecasts use frozen development checkpoints and do not
benefit from additional fitting on the development outcomes.

\paragraph{Dependence fitting.}
At each eligible origin, the pair sampler draws 1,000 uniform-node pairs and,
for each relation type, 1,000 pairs from uniformly sampled eligible groups and
1,000 pairs from groups sampled proportionally to $n_g(n_g-1)$. Groups must
contain at least two eligible entities for pair fitting. This is a composite
training distribution, not an unbiased all-pairs likelihood. Borg has 720,000
fitting and 240,000 validation pairs; BPI has 150,000 and 51,000. Features are
standardized on fitting queries and clipped to $[-10,10]$.

The global and static models use AdamW learning rate 0.02 and no weight decay;
context models use 0.001 and weight decay $10^{-4}$. All use pair batches of
4,096, at most 50 epochs, gradient clipping at 5, and validation patience 8.
Context offsets are bounded by $2\tanh(\cdot)$, with mean-squared shrinkage
coefficient 0.01. Gaussian pair probabilities use 32-point quadrature in double
precision. Validation cycles over the 16 horizon-pair combinations. Model
selection uses validation pair likelihood; the reported final metric remains
the independently specified scenario score.

\paragraph{Numerical checks and hardware.}
Experiments use Python 3.10.14, PyTorch 2.5.1 with CUDA 12.4, NumPy 2.2.6,
SciPy 1.14.1, and two NVIDIA RTX A6000 GPUs. Local tests cover replay and label
semantics, causal history retrieval, monotone trajectories, scoring,
cross-chunk factor sharing, permutation controls, residual resampling, and
frozen-protocol verification. A GPU sampling audit with 200,000 draws checks
marginal preservation, repeated-seed determinism, and shared factors across
entity chunks. These are implementation checks, not evidence that a fitted
distribution is calibrated in the real system.

GPU projection aggregation uses integer counts. Independent CPU and CUDA
random streams need not produce identical draws. A saved temporal-model
inference audit found a maximum CPU--CUDA probability discrepancy of about
0.0021, with much smaller mean error, from floating-point time-feature
arithmetic. Matching-CUDA inference reproduced saved probabilities to below
$2\times10^{-7}$. Final temporal inference therefore uses CUDA. Runtime tables
measure complete scenario sampling plus projection scoring, not isolated GPU
kernel time; other workloads may share the host.

\paragraph{Uncertainty calculation.}
For each model pair, scores are first averaged over seeds at each origin. We
resample paired chronological differences using circular moving blocks, with
10,000 replicates and bootstrap seed 20261014. Borg uses blocks of 4 and 8
origins; BPI uses 14 and 28. We report percentile 95\% intervals only when at
least two blocks fit within the observed span. Overlap of forecast horizons,
drift, and the small number of independent periods remain limitations. No
multiple-comparison-adjusted significance claim is made from exploratory
tables.

\paragraph{Artifact boundaries.}
The source manifests contain URLs, byte counts, and SHA-256 hashes. Run
manifests preserve feature definitions, data hashes, configurations, seeds,
checkpoints, and scores. The final protocol freezes selected checkpoints and
evaluation code before final-period preparation. Original exploratory reports
are retained and are distinguishable from final reports by their period tags.
Tables and figures are generated from saved results. Public raw data are
obtained from the original providers under their original terms; the BPI
archive is not relicensed as Creative Commons by this artifact.

\section{Supplementary quantitative results}
\begin{table*}[t]
\centering\small
\caption{BORG final projection diagnostics. Equal horizon and origin weights; the projection CSV includes all families and modes; the released origin-score tables retain individual seeds.}\label{tab:borgprojection}
\begin{tabular}{lrrrr}
\toprule
Projection/model & CRPS & Cov. & Interval score & Loss$_9$ \\
\midrule
job: Independent & 0.01957 & 0.914 & 0.3073 & 0.0623 \\
job: Global & 0.01542 & 0.971 & 0.1610 & 0.0312 \\
job: Typed static & 0.01451 & 0.978 & 0.1390 & 0.0282 \\
job: Typed context & 0.01370 & 0.983 & 0.1247 & 0.0270 \\
job: Permuted incidence & 0.01729 & 0.943 & 0.2194 & 0.0439 \\
job: S Typed static & 0.01515 & 0.953 & 0.1485 & 0.0387 \\
job: Historical residual & 0.02088 & 0.909 & 0.3564 & 0.0827 \\
machine: Independent & 0.02547 & 0.965 & 0.2153 & 0.0524 \\
machine: Global & 0.02518 & 0.980 & 0.2011 & 0.0507 \\
machine: Typed static & 0.02504 & 0.976 & 0.2004 & 0.0500 \\
machine: Typed context & 0.02498 & 0.978 & 0.1995 & 0.0498 \\
machine: Permuted incidence & 0.02535 & 0.968 & 0.2105 & 0.0516 \\
machine: S Typed static & 0.02505 & 0.974 & 0.2012 & 0.0502 \\
machine: Historical residual & 0.02763 & 0.942 & 0.2571 & 0.0648 \\
population: Independent & 0.05746 & 0.100 & 1.0853 & 0.2785 \\
population: Global & 0.03877 & 0.988 & 0.2912 & 0.0833 \\
population: Typed static & 0.03807 & 0.763 & 0.2867 & 0.0802 \\
population: Typed context & 0.03821 & 0.783 & 0.2902 & 0.0856 \\
population: Permuted incidence & 0.03808 & 0.765 & 0.2873 & 0.0802 \\
population: S Typed static & 0.03919 & 0.655 & 0.3335 & 0.1028 \\
population: Historical residual & 0.04415 & 0.772 & 0.3923 & 0.1045 \\
\bottomrule
\end{tabular}
\end{table*}

\begin{table*}[t]
\centering\small
\caption{BPI final projection diagnostics. Equal horizon and origin weights; the projection CSV includes all families and modes; the released origin-score tables retain individual seeds.}\label{tab:bpiprojection}
\begin{tabular}{lrrrr}
\toprule
Projection/model & CRPS & Cov. & Interval score & Loss$_9$ \\
\midrule
population: Independent & 0.03421 & 0.219 & 0.5710 & 0.0674 \\
population: Global & 0.02728 & 0.869 & 0.2085 & 0.0847 \\
population: Typed static & 0.02857 & 0.668 & 0.2795 & 0.0697 \\
population: Typed context & 0.02824 & 0.720 & 0.2589 & 0.0727 \\
population: Permuted incidence & 0.02852 & 0.671 & 0.2781 & 0.0696 \\
population: S Typed static & 0.03413 & 0.322 & 0.4709 & 0.0659 \\
population: Historical residual & 0.02548 & 0.746 & 0.2476 & 0.0722 \\
handler: Independent & 0.06599 & 0.853 & 0.5660 & 0.1680 \\
handler: Global & 0.06462 & 0.911 & 0.4998 & 0.1733 \\
handler: Typed static & 0.06440 & 0.919 & 0.4920 & 0.1747 \\
handler: Typed context & 0.06436 & 0.935 & 0.4792 & 0.1789 \\
handler: Permuted incidence & 0.06536 & 0.881 & 0.5364 & 0.1692 \\
handler: S Typed static & 0.06640 & 0.892 & 0.5036 & 0.1697 \\
handler: Historical residual & 0.06225 & 0.916 & 0.5025 & 0.1794 \\
\bottomrule
\end{tabular}
\end{table*}

